\documentclass[12pt]{article}
\usepackage[margin=1in]{geometry}
\usepackage{fontspec}
\usepackage{xeCJK}
\usepackage{hyperref}
\usepackage{parskip}
\setmainfont{TeX Gyre Termes}
\setCJKmainfont{Noto Serif CJK SC}
\hypersetup{colorlinks=true, linkcolor=black, urlcolor=blue}

\title{A Bilingual Blog Note / 一篇双语博客小记}
\author{Your Name}
\date{April 2026}

\begin{document}
\maketitle

This template can also be used for bilingual writing.

这份模板也可以直接用来写中英文混合的小博客。

\section*{English}
Sometimes a post is just a compact note, not a full article.

\section*{中文}
有时候博客不需要很长，只要把一个想法说清楚就够了。

\end{document}
